\section{Percolation of spin patterns}

\begin{definition}[Infinite clusters, oceans, Georgii (18.4)--(18.7)]
    \label{def:lo-18-6}
    \lean{SimpleGraph.infiniteClusters, SimpleGraph.IsOceanIn, SimpleGraph.IsOceanIn.infinite, SimpleGraph.IsOceanIn.existsUnique_infinite_supp, SimpleGraph.oceanPart, SimpleGraph.IsOceanIn.subset_oceanPart, SimpleGraph.oceanPart_image, MeasureTheory.GibbsMeasure.Peierls.infinite_of_isOceanIn, MeasureTheory.GibbsMeasure.Peierls.existsUnique_infinite_supp_of_isOceanIn, MeasureTheory.GibbsMeasure.Peierls.oceanPart_nonempty_iff_exists_isOceanIn, MeasureTheory.GibbsMeasure.Peierls.oceanPart_image_add_right, MeasureTheory.GibbsMeasure.Peierls.oceanPart_image_refl}
    \leanok

    For a graph and a set of vertices, $\xi(G,\omega)$ is the union of the infinite clusters of
    the induced subgraph. A set $R$ is an ocean in $P$ if $R \subseteq P$ and $P \setminus R$ has no
    infinite cluster; on $\mathbb Z^2$ every ocean is infinite, a set containing an ocean has a
    unique infinite cluster, itself an ocean, and $\xi^0_P$, the ocean part, is the unique maximal
    ocean, nonempty iff some ocean exists; it is equivariant under the lattice symmetries.
\end{definition}

\begin{definition}[Local ground states and the weight $t(G,\Phi)$, Georgii (18.1)--(18.3), (18.8)]
    \label{def:lo-18-8}
    \uses{def:rp-17-18}
    \lean{MeasureTheory.GibbsMeasure.localGroundEnergy, MeasureTheory.GibbsMeasure.localGroundEnergy_eq_sSup, MeasureTheory.GibbsMeasure.localGroundStates, MeasureTheory.GibbsMeasure.measure_localGroundStates_pos, MeasureTheory.GibbsMeasure.IsRSymmetric, MeasureTheory.GibbsMeasure.isRSymmetric_localGroundStates, MeasureTheory.GibbsMeasure.torusPattern, MeasureTheory.GibbsMeasure.mem_torusPattern_of_isRSymmetric, MeasureTheory.GibbsMeasure.energyOutside, MeasureTheory.GibbsMeasure.groundStateCost, MeasureTheory.GibbsMeasure.patternWeight, MeasureTheory.GibbsMeasure.patternWeight_eq_iInf}
    \leanok

    For a $C$-potential, $m_\Phi$ is the $\lambda^C$-essential infimum of $\Phi_C$ (18.1) and
    $G_\varepsilon(\Phi) = \{\Phi_C \leq m_\Phi + \varepsilon\}$ the set of local ground states
    (18.2), of positive $\lambda^C$-measure. A set $G \subseteq E^C$ is $r$-symmetric if invariant
    under the reflections of the cube, and $V(G,\omega) = \{i : r^i\omega_{C(i)} \in G\}$ is the
    pattern set (18.3), the reflections being immaterial for $r$-symmetric $G$. The weight (18.8) is
    $t(G,\Phi) = e^{-\operatorname{ess\,inf}_{G^c}(\Phi_C - m_\Phi)}\,
    \tilde\lambda^C(G^c)^{1/|C|}\,\inf_{\delta > 0} e^\delta/\tilde\lambda^C(G_\delta(\Phi))$.
\end{definition}

\begin{lemma}[Georgii (18.9)]
    \label{lem:lo-18-9}
    \uses{def:lo-18-8}
    \lean{MeasureTheory.GibbsMeasure.patternWeight_anti, MeasureTheory.GibbsMeasure.groundStateCost_eq_iInf, MeasureTheory.GibbsMeasure.patternWeight_le_of_abs_sub_le, MeasureTheory.GibbsMeasure.patternWeight_smul_le, MeasureTheory.GibbsMeasure.tendsto_patternWeight_smul_atTop}
    \leanok

    $t(\cdot,\Phi)$ is decreasing; the identity (18.9)(2) holds with the essential supremum
    unfolded, as an infimum over pairs $(M, c)$ with $\lambda^C(M) > 0$ and $\Phi_C \leq c$ a.e.\ on
    $M$ (Georgii's $\lambda$-$\sup_M\Phi_C$ is $+\infty$ when $\Phi_C$ is essentially unbounded on
    $M$, where a literal transcription through Mathlib's real essential supremum would be false);
    $t(G,\Psi) \leq e^{4\eta}\,t(G,\Phi)$ when $\|\Psi_C - \Phi_C\| \leq \eta$, which is the
    upper semicontinuity (3); and $t(G_\varepsilon(\Phi),\beta\Phi) \leq
    e^{-\beta\varepsilon/2}/\tilde\lambda^C(G_{\varepsilon/2}(\Phi)) \to 0$ as $\beta \to \infty$.
\end{lemma}
\begin{proof}
    \leanok
\end{proof}

\begin{lemma}[Georgii (18.10), on the torus]
    \label{lem:lo-18-10}
    \uses{def:lo-18-8, cor:rp-17-17-gibbs}
    \lean{MeasureTheory.GibbsMeasure.pi_forall_cubeView_evenSite, MeasureTheory.GibbsMeasure.pow_le_pi_forall_cubeView, MeasureTheory.GibbsMeasure.periodicGibbs_forall_notMem_le, MeasureTheory.GibbsMeasure.le_periodicGibbs_univ, MeasureTheory.GibbsMeasure.periodicGibbsDist_forall_mem_compl_le, MeasureTheory.GibbsMeasure.periodicGibbsDist_pow_le_pow, MeasureTheory.GibbsMeasure.periodicGibbsDist_forall_notMem_torusPattern_le_patternWeight, MeasureTheory.measurePreserving_blocks}
    \leanok

    For $r$-symmetric measurable $G$ and any finite set $D$ of sites of the torus $\Lambda(N)$,
    ${}^\circ\gamma_{\Lambda(N)}^\Phi(D \cap V(G,\cdot) = \emptyset) \leq t(G,\Phi)^{|D|}$. This
    is the estimate for the periodic Gibbs distributions; its passage to the limits
    $\mu \in \mathcal G_0(\Phi)$ waits on the transport of ${}^\circ\gamma_{\Lambda(N)}^\Phi$ to a
    random field on $\mathbb Z^d$.
\end{lemma}
\begin{proof}
    \leanok
    Georgii's proof: the chessboard estimate of Corollary \ref{cor:rp-17-17-gibbs} for the
    indicators $1_{G^c}$, the numerator bounded through the essential infimum of $\Phi_C$ off $G$
    and the product over the even sublattice (which is where the exponent $1/|C|$ comes from), and
    the partition function bounded below by $e^{-|\Lambda|(m_\Phi+\delta)}\lambda^C(G_\delta)^{|\Lambda|}$.
\end{proof}

\begin{lemma}[Georgii (18.13), (18.14)]
    \label{lem:lo-18-14}
    \uses{def:lo-18-6, lem:ising-peierls}
    \lean{MeasureTheory.GibbsMeasure.Peierls.crossingBound, MeasureTheory.GibbsMeasure.Peierls.crossingBound_le_one, MeasureTheory.GibbsMeasure.Peierls.crossingBound_le_four_mul, MeasureTheory.GibbsMeasure.Peierls.le_measure_mem_infiniteClusters, MeasureTheory.GibbsMeasure.Peierls.le_measure_forall_mem_infiniteClusters}
    \leanok

    Let $z(t) = 1 \wedge \sum_{\ell \geq 1}\ell\,5^{\ell-1}t^\ell$, so $z(t) \leq 4t$ for
    $t \leq 1/20$. If a random set $W$ satisfies $\mu(D \cap W = \emptyset) \leq t^{|D|}$ for every
    finite $D$, then the origin lies in an infinite cluster of $W$ within each quadrant with
    probability at least $1 - z(t)$, and in all four at once with probability at least $1 - 4z(t)$.
    No measurability of $W$ is needed: the proof is countable subadditivity over $*$-crossings.
\end{lemma}
\begin{proof}
    \leanok
    A quadrant in which the origin's cluster is finite is cut by a $*$-connected path of sites
    avoiding $W$; there are at most $\ell\,5^{\ell-1}$ such paths of length $\ell$ through a given
    site, and each costs $t^\ell$.
\end{proof}

\begin{remark}[Georgii (18.15)]
    \label{rmk:lo-18-15}

    Poincar\'e recurrence, $\mu(A \setminus A_\infty(\theta)) = 0$ for a measure-preserving
    $\theta$ and a finite $\mu$, is Mathlib's \texttt{MeasurePreserving.conservative} with
    \texttt{Conservative.ae\_mem\_imp\_frequently\_image\_mem}
    (\texttt{Mathlib.Dynamics.Ergodic.Conservative}, not yet in this import graph).
\end{remark}
